\scanpage{19}
\margin{\unsure{\ensuremath{\Vert\;\Vert}}}\margin{Notes anciennes}
$\varphi$ forme symétrique resp. alternée sur un motif $M$ de
poids $p$, $\varphi$ à valeurs dans $Q(p)$, $\varphi$ non
degénérée

$\varphi$ définit un isomorphisme
\[
\mathrm{End}(M)\mathrel{\mathop{\longleftarrow}\limits^{\sim}}\deleted{\unsure{\mathrm{Hom}(M,M^{\vee})}}\Phi(M)
\]
en associant à toute forme $\psi$ l'endomorphisme
$\widetilde\psi^{-1}\widetilde\varphi$ de $M$, et à tout
endomorphisme $u$ de $M$ la
forme $\varphi\circ(u\otimes\mathrm{id}$, \deleted{donc} \unsure{correspondant|correspondant à}

\diagram{\begin{tikzcd}[ampersand replacement=\&]
M \arrow[r,shift left=1.4ex,"\widetilde\varphi"] \arrow[r,shift right=1.4ex,"\widetilde\psi"'] \& M^*=M^{\vee}(p)
\end{tikzcd}}

\margin{$\varphi(x,y)=\deleted{\unsure{\varphi}}(y,\widetilde\varphi(x))$}

$\widetilde\varphi\circ u$. Alors la \uline{symétrie} de $\Phi(M)$ devient
une application de $\mathrm{End}(M)$ en lui-$=$, notée $u\leadsto$
$u^\varphi$, qui n'est autre que l'\uline{adjoint}
relativement à $\varphi$,
% archive sheet number: 18
